% ch03_u06 —— 原书 p.112–p.121（3.7 末的无编号段落、3.8 全部 3.8.1–3.8.23）
% p.112 首段无段号，接 p.111 末 3.7.17 证明之后；3.8.23 在 p.122 续，不属本单元。
\SourcePage{112}

根据该引理，映射（124）的核 $\mathfrak{a}_{\mathfrak{F}}$ 作为李代数与 $(\mathfrak{g}_{\mathrm{univ}})_{\mathfrak{F}}$ 一致，后者是 3.5.11 中李代数 $\mathfrak{g}_{\mathrm{univ}}$ 在 $\mathfrak{F}$ 处的纤维。映射（124）=（125）是满射，因为 $\nabla:\mathfrak{g}^{O}_{\mathfrak{F}}\to\mathfrak{g}^{O}_{\mathfrak{F}}\otimes\omega_{O}$ 是满射。容易证明，（121）与（122）是同胚，且映射（124）是开映射。

类似地可以证明，态射（123）是满射且开的，它的核可典范地等同于带离散拓扑的 $\mathfrak{g}_{\mathrm{univ}}$（这一等同对每个 $\mathfrak{F}\in\operatorname{Op}_{\mathfrak{g}}(O)$ 诱导出上述同构 $\mathfrak{a}_{\mathfrak{F}}\xrightarrow{\sim}(\mathfrak{g}_{\mathrm{univ}})_{\mathfrak{F}}$）。引理 3.5.12 给出连续的李代数胚态射 $f:I/I^{2}\to\mathfrak{a}_{\mathfrak{g}}$，使得图表
\[
\begin{tikzcd}
  0 \arrow[r] & \mathfrak{g}_{\mathrm{univ}} \arrow[r] \arrow[d, "\mathrm{id}"'] & I/I^{2} \arrow[r] \arrow[d, "f"']
    & \operatorname{Der}A_{\mathfrak{g}}(O) \arrow[r] \arrow[d, "\mathrm{id}"] & 0 \\
  0 \arrow[r] & \mathfrak{g}_{\mathrm{univ}} \arrow[r] & \mathfrak{a}_{\mathfrak{g}} \arrow[r]
    & \operatorname{Der}A_{\mathfrak{g}}(O) \arrow[r] & 0
\end{tikzcd}
\]
交换。由于图表的两行在拓扑意义下正合，$f$ 是拓扑同构。显然 $f$ 是 $\operatorname{Aut}O$-等变的。

\bdsec{3.8. oper 的奇点}

\bdp{3.8.1} 设 $U$ 是曲线 $X$ 的开稠子集。我们要把归纳概形 $\operatorname{Op}_{\mathfrak{g}}(U)$ 表示成若干闭子概形 $\operatorname{Op}_{\mathfrak{g},D}(X)$ 的并，其中 $D$ 跑遍 $X$ 的有限子概形，且满足 $D\cap U=\emptyset$。

由 3.1.9，有典范同构 $\underline{\operatorname{Op}}_{\mathfrak{g}}(U)\xrightarrow{\sim}\operatorname{Op}_{\mathfrak{g}}(U)$。这里 $\underline{\operatorname{Op}}_{\mathfrak{g}}(U)$ 是由 $\Gamma(U,\omega_{X}^{\otimes 2})$-torsor $\operatorname{Op}_{\mathfrak{sl}_{2}}(U)$ 经某个嵌入 $\Gamma(U,\omega_{X}^{\otimes 2})\subset\Gamma(U,V_{\omega_{X}})$ 诱导出的 $\Gamma(U,V_{\omega_{X}})$-torsor。该嵌入与 $V=V_{\mathfrak{g}}$ 的定义见 3.1.9。提醒一句：$V$ 是带 $\mathbb{G}_{m}$-作用（即分次）的向量空间，$V_{\omega_{X}}$ 表示用 $\mathbb{G}_{m}$-torsor $\omega_{X}$ 扭 $V$ 所得的扭。有 $\dim V=r:=\operatorname{rank}\mathfrak{g}$，且 $V$ 各分次分量的次数等于 $\mathfrak{g}$ 上「基本」不变多项式的次数 $d_{1},\ldots,d_{r}$。

\SourcePage{113}

若 $D$ 是 $X$ 的有限子概形，则有典范嵌入 $V_{\omega_{X}}\hookrightarrow V_{\omega_{X}(D)}$。记 $\overline{\operatorname{Op}}_{\mathfrak{g},D}(X)$ 为由 $\Gamma(X,V_{\omega_{X}})$-torsor $\underline{\operatorname{Op}}_{\mathfrak{g}}(X)$ 诱导出的 $\Gamma(X,V_{\omega_{X}(D)})$-torsor。显然 $\overline{\operatorname{Op}}_{\mathfrak{g},D}(X)$ 是 $\underline{\operatorname{Op}}_{\mathfrak{g}}(X\setminus D)$ 的闭子概形。记 $\operatorname{Op}_{\mathfrak{g},D}(X)$ 为 $\overline{\operatorname{Op}}_{\mathfrak{g},D}(X)$ 在 $\operatorname{Op}_{\mathfrak{g}}(X\setminus D)$ 中的像。若 $D\subset D'$，则 $\operatorname{Op}_{\mathfrak{g},D}(X)\subset\operatorname{Op}_{\mathfrak{g},D'}(X)$。对任意开稠子集 $U\subset X$，有
\[
\operatorname{Op}_{\mathfrak{g}}(U)=\bigcup_{D\cap U=\emptyset}\operatorname{Op}_{\mathfrak{g},D}(X).
\]

3.8.23 中将给出 $\operatorname{Op}_{\mathfrak{g},D}(X)$ 的「内蕴」描述，它不使用同构 $\underline{\operatorname{Op}}_{\mathfrak{g}}\xrightarrow{\sim}\operatorname{Op}_{\mathfrak{g}}$。该描述的局部版本见 3.8.7--3.8.10。

\bdp{3.8.2} 现在可以给出 2.8.6 中问题的答案：
\begin{equation}\tag{128}
  N_{\Delta}(G)=\operatorname{Op}_{{}^{L}\mathfrak{g},\Delta}(X).
\end{equation}
$N_{\Delta}(G)$ 定义为归纳概形 $N'_{\Delta}(G)$ 的子概形，后者经 Feigin--Frenkel 同构典范地等同于 $\operatorname{Op}_{{}^{L}\mathfrak{g}}(X\setminus\Delta)$。（128）是 $\operatorname{Op}_{{}^{L}\mathfrak{g}}(X\setminus\Delta)$ 的子概形之间的等式。

我们不证明（128）。3.8.6 中将给出一个提示。

\bdp{3.8.3} 3.8.1 中 $\operatorname{Op}_{\mathfrak{g},D}(X)$ 的定义在如下局部情形有意义：
\[
X=\operatorname{Spec}O,\qquad O:=\mathbb{C}[[t]],\qquad D=\operatorname{Spec}O/t^{n}O.
\]
此时记 $\operatorname{Op}_{\mathfrak{g},n}(O)$ 以代替 $\operatorname{Op}_{\mathfrak{g},D}(X)$。$\operatorname{Op}_{\mathfrak{g},n}(O)$ 是归纳概形 $\operatorname{Op}_{\mathfrak{g}}(K)$ 的闭子概形。当然 $\operatorname{Op}_{\mathfrak{g},0}(O)=\operatorname{Op}_{\mathfrak{g}}(O)$，$\operatorname{Op}_{\mathfrak{g},n}(O)\subset\operatorname{Op}_{\mathfrak{g},n+1}(O)$，且 $\operatorname{Op}_{\mathfrak{g}}(K)$ 是诸 $\operatorname{Op}_{\mathfrak{g},n}(O)$ 的归纳极限。

由 3.7.12，$A_{\mathfrak{g}}(K)$ 是 $\operatorname{Op}_{\mathfrak{g}}(K)$ 上正则函数构成的代数。记 $I_{n}$ 为 $A_{\mathfrak{g}}(K)$ 中对应于 $\operatorname{Op}_{\mathfrak{g},n}(O)\subset\operatorname{Op}_{\mathfrak{g}}(K)$ 的理想。显然 $I_{n}\supset I_{n+1}$，且 $I_{0}$ 就是 3.7.16(i) 中的理想 $I$。诸理想 $I_{n}$ 构成 $A_{\mathfrak{g}}(K)$ 中 $0$ 的邻域基。

\bdp{3.8.4} 下面给出 $A_{\mathfrak{g}}(K)$ 与 $I_{n}$ 的显式描述。沿用 3.5.6 的记号，于是 $\operatorname{Spec}K$ 上的 $\mathfrak{g}$-oper 与满足 $u_{j}(t)\in\mathbb{C}((t))$ 的算子（64）一一对应。把 $u_{j}(t)$ 写成 $\sum_{k}\tilde{u}_{jk}t^{k}$。则 $A_{\mathfrak{g}}(K)=\mathbb{C}[[\ldots,\tilde{u}_{j,-1},\tilde{u}_{j0},\tilde{u}_{j1},\ldots]]$（用记号（98））。理想 $I_{n}$ 由诸 $\tilde{u}_{jk}$（$k<-d_{j}n$）拓扑生成。\SourcePage{114} 3.5.6 中的 $u_{jk}$ 是 $\tilde{u}_{jk}$ 在 $A_{\mathfrak{g}}(O)=A_{\mathfrak{g}}(K)/I$ 中的像。

容易描述 $\operatorname{Der}K$ 在 $A_{\mathfrak{g}}(K)$ 上的作用。特别地
\begin{equation}\tag{129}
  L_{0}\tilde{u}_{jk}=(d_{j}+k)\tilde{u}_{jk}.
\end{equation}

与整体情形一样（见 3.1.12--3.1.14），$\operatorname{Op}_{\mathfrak{g}}(K)$ 的坐标环 $A_{\mathfrak{g}}(K)$ 带典范滤链。其第 $i$ 项由 $\tilde{u}_{jk}$ 的「多项式」组成，这些多项式的加权次数 $\le i$；这里约定 $\tilde{u}_{jk}$ 的权为 $d_{j}$。

\bdp{3.8.5} \emph{命题.} 理想 $I_{n}\subset A_{\mathfrak{g}}(K)$ 由诸空间 $A_{i}^{m}$（$m<i(1-n)$）拓扑生成，其中 $A_{i}^{m}$ 是 $A_{\mathfrak{g}}(K)$ 的滤链第 $i$ 项中满足 $L_{0}a=ma$ 的元素 $a$ 的集合。\hfill$\Box$

同构 $A_{{}^{L}\mathfrak{g}}(K)\xrightarrow{\sim}\mathfrak{z}$（见（113））保持滤链，且是 $\operatorname{Aut}K$-等变的。于是命题 3.8.5 推出如下陈述。

\bdp{3.8.6} \emph{命题.} Feigin--Frenkel 同构 $A_{{}^{L}\mathfrak{g}}(K)\xrightarrow{\sim}\mathfrak{z}$ 把 $I_{n}\subset A_{{}^{L}\mathfrak{g}}(K)$ 映到 3.7.9 中的理想 $I_{n}$。

这是（128）的证明中用到的组成部分之一。

\bdp{3.8.7} 我们要用「自然」的语言（不使用同构（43））描述 $\operatorname{Op}_{\mathfrak{g},n}(O)$。记 $\mathfrak{g}^{+}$ 为 $\mathfrak{g}$ 的如下局部闭既约子概形：它由所有满足下述条件的 $a\in\mathfrak{g}$ 组成——对正根 $\alpha$，若 $\alpha$ 非单则 $a_{-\alpha}=0$，若 $\alpha$ 单则 $a_{-\alpha}\ne 0$（$a_{-\alpha}$ 是 $a$ 在根子空间 $\mathfrak{g}^{-\alpha}$ 中的分量）。于是对任意 $\mathbb{C}$-代数 $R$，集合 $\mathfrak{g}^{+}(R)$ 由 $a\in\mathfrak{g}\otimes R$ 组成，使得对每个非单的 $\alpha>0$ 有 $a_{-\alpha}=0$，且对每个单的 $\alpha$，$a_{-\alpha}$ 生成 $R$-模 $\mathfrak{g}^{-\alpha}\otimes R$。

回忆：$\operatorname{Spec}K$ 上的 $\mathfrak{g}$-oper 是算子 $\frac{d}{dt}+q(t)$（$q\in\mathfrak{g}^{+}(K)$）的 $B(K)$-共轭类。这里 $B$ 是 $\mathfrak{g}$ 所对应的伴随群 $G$ 的 Borel 子群。

\SourcePage{115}

\bdp{3.8.8} \emph{定义.} $\operatorname{Spec}O$ 上的 $(\le n)$-奇异 $\mathfrak{g}$-oper 是算子 $\frac{d}{dt}+t^{-n}q(t)$（$q\in\mathfrak{g}^{+}(O)$）的 $B(O)$-共轭类。

\emph{注.}
\begin{enumerate}[label=(\roman*),leftmargin=2.2em]
\item $B(O)$ 在算子集 $\{\frac{d}{dt}+t^{-n}q(t)|q\in\mathfrak{g}^{+}(O)\}$ 上的作用是自由的。事实上，$B(K)$ 在 $\{\frac{d}{dt}+q(t)|q\in\mathfrak{g}^{+}(K)\}$ 上的作用是自由的（见 3.1.4）。
\item $n=0$ 时得到 $\operatorname{Spec}O$ 上 $\mathfrak{g}$-oper 的通常概念。
\end{enumerate}

\bdp{3.8.9} \emph{命题.} 映射「$\operatorname{Spec}O$ 上的 $(\le n)$-奇异 $\mathfrak{g}$-oper」$\to\operatorname{Op}_{\mathfrak{g}}(K)$ 是单射。其像等于 $\operatorname{Op}_{\mathfrak{g},n}(O)$。

\emph{证明.} 沿用 3.5.6 的记号。对每个 $v_{1},\ldots,v_{r}\in\mathbb{C}[[t]]$，算子
\begin{equation}\tag{130}
  \frac{d}{dt}+t^{-n}(i(f)+v_{1}(t)e_{1}+\cdots+v_{r}(t)e_{r})
\end{equation}
定义了 $\operatorname{Spec}O$ 上一个 $(\le n)$-奇异 $\mathfrak{g}$-oper。容易证明，这在算子（130）与 $\operatorname{Spec}O$ 上的 $(\le n)$-奇异 $\mathfrak{g}$-oper 之间给出双射。现在用 $B(K)$-共轭把（130）化成「典范形式」（64）。用 $t^{-n\check\rho}$ 共轭（130），得到
\begin{equation}\tag{131}
  \frac{d}{dt}+i(f)+n\check\rho t^{-1}+t^{-nd_{1}}v_{1}(t)e_{1}+\cdots+t^{-nd_{r}}v_{r}(t)e_{r}.
\end{equation}
% [Tutor Note: 原书 p.115 把 $\exp(-ne_{1}/2t)$ 与 $n(n-2)/4t^{2}$ 印作线性分数，
%  字面亦可读作 $(ne_{1}/2)t$、$(n(n-2)/4)t^{2}$；由本段共轭计算（$n=1$ 时应得 $-1/4t^{2}$）
%  可知应为 $e_{1}/(2t)$、$n(n-2)/(4t^{2})$。底本如此，正文照录。]
为消去 $n\check\rho t^{-1}$，再用 $\exp(-ne_{1}/2t)$ 共轭（131），得到算子（64），其中
\[
\begin{aligned}
  u_{j}(t)&=t^{-nd_{j}}v_{j}(t) \qquad \text{对 } j>1,\\
  u_{1}(t)&=t^{-nd_{1}}v_{1}(t)+n(n-2)/4t^{2},\qquad d_{1}=2.
\end{aligned}
\]
显然 $v_{j}\in\mathbb{C}[[t]]$ 当且仅当 $u_{j}\in t^{-nd_{j}}\mathbb{C}[[t]]$。\hfill$\Box$

\bdp{3.8.10} 若把 $\operatorname{Op}_{\mathfrak{g},n}(O)$ 的点看作 $\operatorname{Spec}O$ 上的 $(\le n)$-奇异 $\mathfrak{g}$-oper，则典范嵌入 $\operatorname{Op}_{\mathfrak{g},n}(O)\hookrightarrow\operatorname{Op}_{\mathfrak{g},n+1}(O)$ 把 $\frac{d}{dt}+t^{-n}q(t)$（$q\in\mathfrak{g}^{+}(O)$）的 $B(O)$-共轭类映到 $t^{\check\rho}(\frac{d}{dt}+t^{-n}q(t))t^{-\check\rho}$ 的 $B(O)$-共轭类（这是良定义的，因为 $t^{\check\rho}B(O)t^{-\check\rho}\subset B(O)$）。

\SourcePage{116}

\bdp{3.8.11} 记 $\operatorname{Inv}(\mathfrak{g})$ 为 $\mathfrak{g}$ 上 $G$-不变多项式构成的代数。有典范态射 $\mathfrak{g}\to\operatorname{Spec}\operatorname{Inv}(\mathfrak{g})=W\backslash\mathfrak{h}$，其中 $W$ 是 Weyl 群。

设有 $\operatorname{Spec}O$ 上的 $(\le 1)$-奇异 $\mathfrak{g}$-oper，即算子 $\frac{d}{dt}+t^{-1}q(t)$（$q\in\mathfrak{g}^{+}(O)$）的 $B(O)$-共轭类。$q(0)\in\mathfrak{g}$ 在 $\operatorname{Spec}\operatorname{Inv}(\mathfrak{g})$ 中的像称为该 oper 的\emph{留数}。于是我们定义了\emph{留数映射}
\begin{equation}\tag{132}
  \operatorname{Res}:\operatorname{Op}_{\mathfrak{g},1}(O)\to\operatorname{Spec}\operatorname{Inv}(\mathfrak{g})=W\backslash\mathfrak{h}.
\end{equation}
它是满射。因此它诱导嵌入
\begin{equation}\tag{133}
  \operatorname{Inv}(\mathfrak{g})\hookrightarrow A_{\mathfrak{g}}(K)/I_{1}
\end{equation}
（回忆：$A_{\mathfrak{g}}(K)/I_{1}$ 是 $\operatorname{Op}_{\mathfrak{g},1}(O)$ 的坐标环；见 3.8.3）。

\bdp{3.8.12} \emph{命题.} $\operatorname{Res}(\operatorname{Op}_{\mathfrak{g}}(O))\subset W\backslash\mathfrak{h}$ 由单点组成，该点是 $-\check\rho\in\mathfrak{h}$ 的像。

\emph{注.} 我们宁愿忘记 $-\check\rho$ 与 $\check\rho$ 在 $W\backslash\mathfrak{h}$ 中有相同的像。

\emph{证明.} 需要计算 3.8.10 所述映射 $\operatorname{Op}_{\mathfrak{g}}(O)\to\operatorname{Op}_{\mathfrak{g},1}(O)$ 与映射（132）的复合。若 $q(t)\in\mathfrak{g}^{+}(O)$，则 $t^{\check\rho}(\frac{d}{dt}+q(t))t^{-\check\rho}=\frac{d}{dt}+\frac{a-\check\rho}{t}+\{\text{正则项}\}$，其中 $a$ 属于对应于单负根的根子空间之和。现在 $a-\check\rho$ 与 $-\check\rho$ 在 $W\backslash\mathfrak{h}$ 中有相同的像。\hfill$\Box$

\bdp{3.8.13} \emph{命题.} 设 $f\in A_{\mathfrak{g}}(K)/I_{1}$，即 $f$ 是 $\operatorname{Op}_{\mathfrak{g},1}(O)$ 上的正则函数。则以下条件等价：
\begin{enumerate}[label=(\roman*),leftmargin=2.2em]
\item $f\in\operatorname{Inv}(\mathfrak{g})$，其中 $\operatorname{Inv}(\mathfrak{g})$ 按（133）与其像等同；
\item $f$ 是 $\operatorname{Aut}^{0}O$-不变的；
\item $L_{0}f=0$。
\end{enumerate}

\emph{证明.} 显然 (i)$\Rightarrow$(ii)$\Rightarrow$(iii)。下面从 (ii) 推出 (i)。考虑 $\operatorname{Spec}O$ 上的一个 $(\le 1)$-奇异 $\mathfrak{g}$-oper。它是联络 $\frac{d}{dt}+t^{-1}q(t)$（$q\in\mathfrak{g}^{+}(O)$）的 $B(O)$-共轭类。若把 $t$ 换成 $\lambda t$，该联络换成 $\frac{d}{dt}+t^{-1}q(\lambda t)$。由于 $L_{0}f=0$，$f$ 在该联络上的值\SourcePage{117}不依赖于 $\lambda$，因而只依赖于 $q(0)\in\mathfrak{g}^{+}$（因为 $\lim_{\lambda\to 0}q(\lambda t)=q(0)$）。剩下的只需用如下事实：$\mathfrak{g}^{+}$ 上 $B$-不变的正则函数可延拓为 $\mathfrak{g}$ 上 $G$-不变的多项式（见 [Ko63] 的定理 0.10）。\hfill$\Box$

\bdp{3.8.14} \emph{注.} 由 3.8.4，代数 $A_{\mathfrak{g}}(K)/I_{1}$ 由 $\tilde{u}_{jk}$（$k\ge -d_{j}$）自由生成，其中 $\tilde{u}_{jk}\in A_{\mathfrak{g}}(K)/I_{1}$ 是 $\tilde{u}_{jk}\in A_{\mathfrak{g}}(K)$ 的像。由 3.8.13 与（129），$\operatorname{Inv}(\mathfrak{g})\subset A_{\mathfrak{g}}(K)/I_{1}$ 由 $v_{j}:=\tilde{u}_{j,-d_{j}}$ 生成。同构 $\operatorname{Spec}\mathbb{C}[v_{1},\ldots,v_{r}]\xrightarrow{\sim}\operatorname{Spec}\operatorname{Inv}(\mathfrak{g})$ 是复合 $\operatorname{Spec}\mathbb{C}[v_{1},\ldots,v_{r}]\to\mathfrak{g}\to\operatorname{Spec}\operatorname{Inv}(\mathfrak{g})$，其中第一个映射等于 $i(f)-\check\rho+v_{1}e_{1}+\cdots+v_{r}e_{r}$（用 3.5.6 的记号）。

\bdp{3.8.15} 我们将证明定理 3.6.11。在 3.8.16--3.8.17 中，将叙述 Feigin--Frenkel 同构（113）的一个性质。该性质把定理 3.6.11 归结为某个陈述（见 3.8.19），后者只涉及 oper 与 Gelfand--Dikii 括号。该陈述将在 3.8.20--3.8.22 中证明。

\bdp{3.8.16} 沿用 3.5.17 的记号。此外，若 $\operatorname{Der}O$ 作用在向量空间 $V$ 上，则令 $V^{0}:=\{v\in V|L_{0}v=0\}$。

如 3.6.9 所解释，3.6.8 中的映射 $\pi$ 诱导态射
\begin{equation}\tag{134}
  (\mathfrak{z}/\mathfrak{z}\cdot\mathfrak{z}^{<0})^{0}=(\mathfrak{z}/\mathfrak{z}\cdot\mathfrak{z}^{<0})^{\le 0}
  =\mathfrak{z}^{\le 0}/(\mathfrak{z}\cdot\mathfrak{z}^{<0}\cap\mathfrak{z}^{\le 0})\to C
\end{equation}
其中 $C$ 是 $U\mathfrak{g}$ 的中心。现在（113）诱导同构
\begin{equation}\tag{135}
  (\mathfrak{z}/\mathfrak{z}\cdot\mathfrak{z}^{<0})^{0}\xrightarrow{\sim}(A_{{}^{L}\mathfrak{g}}(K)/I_{1})^{0},
\end{equation}
因为由 3.8.5，$I_{1}=A_{{}^{L}\mathfrak{g}}(K)\cdot A_{{}^{L}\mathfrak{g}}(K)^{<0}$。由 3.8.13，（135）的右端等于 $\operatorname{Inv}({}^{L}\mathfrak{g})$。于是（134）与（135）给出态射
\begin{equation}\tag{136}
  \operatorname{Inv}({}^{L}\mathfrak{g})\to C.
\end{equation}

记 $\operatorname{Inv}(\mathfrak{h}^{*})$ 为 $\mathfrak{h}^{*}$ 上 $W$-不变多项式构成的代数。由于 ${}^{L}\mathfrak{h}=\mathfrak{h}^{*}$，有典范同构 $\operatorname{Inv}({}^{L}\mathfrak{g})\xrightarrow{\sim}\operatorname{Inv}(\mathfrak{h}^{*})$。此外还有\SourcePage{118} Harish-Chandra 同构 $C\xrightarrow{\sim}\operatorname{Inv}(\mathfrak{h}^{*})$。于是（136）可看作映射
\begin{equation}\tag{137}
  \operatorname{Inv}(\mathfrak{h}^{*})\to\operatorname{Inv}(\mathfrak{h}^{*}).
\end{equation}

\bdp{3.8.17} \emph{定理.}（E. Frenkel，私人通信）

态射（137）把 $f\in\operatorname{Inv}(\mathfrak{h}^{*})$ 映到 $f^{-}$，其中 $f^{-}(\lambda):=f(-\lambda)$，$\lambda\in\mathfrak{h}^{*}$。\hfill$\Box$

\bdp{3.8.18} 利用 3.8.17，可以把图表（84）神秘的左下角换成其 oper 版本。下面的图表（143）本质上就是这样得到的。我们来定义（143）的下方箭头，它是 3.6.9--3.6.10 中所构造映射（83）的 oper 版本。

由 3.8.5，
\begin{equation}\tag{138}
  I_{1}=A_{\mathfrak{g}}(K)\cdot A_{\mathfrak{g}}(K)^{<0}.
\end{equation}

由 3.8.13，有典范同构
\begin{equation}\tag{139}
  (A_{\mathfrak{g}}(K)/I_{1})^{0}\xrightarrow{\sim}\operatorname{Inv}(\mathfrak{g}).
\end{equation}

对 $h\in\mathfrak{h}$，记 $m_{h}$ 为 $\operatorname{Inv}(\mathfrak{g})$ 中由在 $h$ 处取零值的多项式构成的极大理想。令 $m:=m_{-\check\rho}$。由 3.8.12，同构（139）诱导
\begin{equation}\tag{140}
  (I/I_{1})^{0}\xrightarrow{\sim}m.
\end{equation}

现在得到
\begin{equation}\tag{141}
  (I/(I^{2}+I_{1}))^{0}\xrightarrow{\sim}m/m^{2}
\end{equation}
（为从（140）得到（141），我们用到
\[
(I^{2})^{0}\subset(I^{0})^{2}+I\cdot I^{<0}\subset(I^{0})^{2}+A_{\mathfrak{g}}(K)\cdot A_{\mathfrak{g}}(K)^{<0}=(I^{0})^{2}+I_{1};
\]
见（138））。

对正则的 $h\in\mathfrak{h}$，我们把 $m_{h}/m_{h}^{2}$ 与 $\mathfrak{h}^{*}$ 等同：把 $\mathfrak{h}$ 上的 $W$-不变多项式对应到它在 $h$ 处的微分。特别地，对 $m=m_{-\check\rho}$，\SourcePage{119}有 $m/m^{2}\xrightarrow{\sim}\mathfrak{h}^{*}$（顺便说一句，若把 $m$ 写成 $m_{\check\rho}$，就会得到另一个同构 $m/m^{2}\xrightarrow{\sim}\mathfrak{h}^{*}$）。

最后，利用（138）把（141）的左端改写成用 $I/I^{2}$ 表达的式子，得到同构
\begin{equation}\tag{142}
  (I/I^{2})^{\le 0}/(A_{\mathfrak{g}}(O)\cdot(I/I^{2})^{<0}\cap(I/I^{2})^{\le 0})\xrightarrow{\sim}\mathfrak{h}^{*}.
\end{equation}

\bdp{3.8.19} \emph{命题.} 图表
\begin{equation}\tag{143}
\begin{tikzcd}[row sep=1.5em]
  (\mathfrak{a}_{\mathfrak{g}})^{\le 0}/(A_{\mathfrak{g}}(O)\cdot\mathfrak{a}_{\mathfrak{g}}^{<0}\cap(\mathfrak{a}_{\mathfrak{g}})^{\le 0})
    \arrow[r, "\sim"] & \mathfrak{h} \arrow[d, "\wr"] \\
  (I/I^{2})^{\le 0}/(A_{\mathfrak{g}}(O)\cdot(I/I^{2})^{<0}\cap(I/I^{2})^{\le 0})
    \arrow[r, "\sim"] \arrow[u, "\wr"] & \mathfrak{h}^{*}
\end{tikzcd}
\end{equation}
交换。这里下方箭头是同构（142），上方箭头是同构（78），左边箭头由同构（120）诱导（后者来自 $A_{\mathfrak{g}}(K)$ 上的 Gelfand--Dikii 括号），右边箭头由定义 Gelfand--Dikii 括号时所用 $\mathfrak{g}$ 上的不变数量积诱导。

该命题将在 3.8.20--3.8.22 中证明。

定理 3.6.11 由 3.8.17 与 3.8.19 推出。（143）的交换性蕴含（84）的反交换性，因为如下图表反交换：
\[
\begin{tikzcd}
  m_{\check\rho}/(m_{\check\rho})^{2} \arrow[rr, "\sim"] \arrow[dr] & & m_{-\check\rho}/(m_{-\check\rho})^{2} \arrow[dl] \\
  & \mathfrak{h}^{*} &
\end{tikzcd}
\]
这里上方箭头由 3.8.17 中的映射 $f\mapsto f^{-}$ 诱导。

\bdp{3.8.20} 我们要叙述命题 3.8.19 的证明中用到的引理。考虑复合
\begin{equation}\tag{144}
  I/I^{2}\to I/(I^{2}+I_{1})\xrightarrow{\sim}\mathfrak{a}_{\mathfrak{g}}/A_{\mathfrak{g}}(O)\cdot\mathfrak{a}_{\mathfrak{g}}^{<0}
  =\mathfrak{a}_{\mathfrak{g}}/\mathfrak{a}_{\mathfrak{n}}=\mathfrak{g}_{\mathrm{univ}}/\mathfrak{n}_{\mathrm{univ}}.
\end{equation}
这里第二个箭头来自（120）与（138）；$\mathfrak{a}_{\mathfrak{n}}$ 与 $\mathfrak{n}_{\mathrm{univ}}$ 已在 3.5.16 中定义，$\mathfrak{a}_{\mathfrak{g}}$ 已在 3.5.11 中定义；等式 $\mathfrak{a}_{\mathfrak{n}}=A_{\mathfrak{g}}(O)\cdot\mathfrak{a}_{\mathfrak{g}}^{<0}$\SourcePage{120} 已在 3.5.18 中证明。$I/I^{2}$ 在 $\mathfrak{F}=(\mathfrak{F}_{B},\nabla)\in\operatorname{Op}_{\mathfrak{g}}(O)$ 上的纤维等于 $\{u\in\mathfrak{g}^{O}_{\mathfrak{F}}|\nabla(u)\in\mathfrak{b}^{O}_{\mathfrak{F}}\otimes\omega_{O}\}/\mathfrak{n}^{O}_{\mathfrak{F}}$（见（122）），而 $\mathfrak{g}_{\mathrm{univ}}/\mathfrak{n}_{\mathrm{univ}}$ 在 $\mathfrak{F}$ 上的纤维等于 $(\mathfrak{g}_{\mathfrak{F}}/\mathfrak{n}_{\mathfrak{F}})_{0}$，即 $\mathfrak{g}_{\mathfrak{F}}/\mathfrak{n}_{\mathfrak{F}}$ 在原点 $0\in\operatorname{Spec}O$ 处的纤维。考虑映射
\[
\varphi,\psi:\{u\in\mathfrak{g}^{O}_{\mathfrak{F}}|\nabla(u)\in\mathfrak{b}^{O}_{\mathfrak{F}}\otimes\omega_{O}\}/\mathfrak{n}^{O}_{\mathfrak{F}}
  \to(\mathfrak{g}_{\mathfrak{F}}/\mathfrak{n}_{\mathfrak{F}})_{0}
\]
其中 $\varphi$ 由（144）诱导，$\psi$ 是在 $0$ 处取值。

\bdp{3.8.21} \emph{引理.} $\varphi=\psi$。

\emph{证明.} 由 3.7.17，$\varphi$ 与 $\psi$ 在 $\mathfrak{a}_{\mathfrak{F}}:=\{u\in\mathfrak{g}^{O}_{\mathfrak{F}}|\nabla(u)=0\}$ 上的限制相等。因此只需证明 $\operatorname{Ker}\varphi\subset\operatorname{Ker}\psi$。显然 $\operatorname{Ker}\varphi=T^{\perp}_{\mathfrak{F}}\operatorname{Op}_{\mathfrak{g},1}(O)$，即 $\operatorname{Op}_{\mathfrak{g},1}(O)$ 在 $\mathfrak{F}$ 处的余法空间。对任意 $q\in\mathfrak{b}^{O}_{\mathfrak{F}}$，oper $\mathfrak{F}_{q}:=(\mathfrak{F}_{B},\nabla+q\cdot\frac{d}{dt})$ 是 $(\le 1)$-奇异的。于是 $\mathfrak{b}^{O}_{\mathfrak{F}}\otimes t^{-1}\omega_{O}$ 在（117）右端的像含于切空间 $T_{\mathfrak{F}}\operatorname{Op}_{\mathfrak{g},1}(O)$ 中。因此 $T^{\perp}_{\mathfrak{F}}\operatorname{Op}_{\mathfrak{g},1}(O)\subset\operatorname{Ker}\psi$。\hfill$\Box$

\bdp{3.8.22} 现在证明 3.8.19。由于（142）的左端等于（141）的左端，可以把 3.8.19 改述如下。

设 $f\in\operatorname{Inv}(\mathfrak{g})$，$f(-\check\rho)=0$。把 $f$ 看作 $A_{\mathfrak{g}}(K)/I_{1}$ 的元素（见（133））。由 3.8.12，$f\in I/I_{1}$。$f$ 在 $I/(I^{2}+I_{1})$ 中的像可看作元素 $\nu\in\mathfrak{g}_{\mathrm{univ}}/\mathfrak{n}_{\mathrm{univ}}$（见（144））。另一方面，设 $\lambda\in\mathfrak{h}^{*}$ 是 $f\in\operatorname{Inv}(\mathfrak{g})$ 在 $\mathfrak{h}$ 上的限制在 $-\check\rho$ 处的微分。为证明 3.8.19，必须证明 $\nu$ 等于 $\lambda$ 在复合
\[
\mathfrak{h}^{*}\xrightarrow{\sim}\mathfrak{h}\subset\mathfrak{h}\otimes A_{\mathfrak{g}}(O)
  =\mathfrak{h}_{\mathrm{univ}}/\mathfrak{n}_{\mathrm{univ}}\subset\mathfrak{g}_{\mathrm{univ}}/\mathfrak{n}_{\mathrm{univ}}
\]
下的像。

由 3.8.21，这等价于如下陈述：设 $\mathfrak{F}=(\mathfrak{F}_{B},\nabla)\in\operatorname{Op}_{\mathfrak{g}}(O)$，$q\in\mathfrak{b}^{O}_{\mathfrak{F}}$，$\mathfrak{F}_{\varepsilon q}:=(\mathfrak{F}_{B},\nabla+\varepsilon q\frac{d}{dt})$，则
\begin{equation}\tag{145}
  \frac{d}{d\varepsilon}f(\operatorname{Res}(\mathfrak{F}_{\varepsilon q}))|_{\varepsilon=0}=\lambda(q_{\mathfrak{h}}(0)),
\end{equation}
其中 $q_{\mathfrak{h}}(t)\in\mathfrak{h}[[t]]$ 是 $q$ 在 $\mathfrak{b}^{O}_{\mathfrak{F}}/\mathfrak{n}^{O}_{\mathfrak{F}}=\mathfrak{h}\otimes O$ 中的像。与 3.8.12 的证明一样可以证明，$\operatorname{Res}(\mathfrak{F}_{\varepsilon q})$ 等于 $-\check\rho+\varepsilon q_{\mathfrak{h}}(0)$ 在 $W\backslash\mathfrak{h}$ 中的像。于是（145）是显然的。

\SourcePage{121}

\bdp{3.8.23} 本小节（读者完全可以跳过）给出 3.8.1 中概形 $\operatorname{Op}_{\mathfrak{g},D}(X)$ 的「内蕴」描述。它由对 3.8.7--3.8.10 作直接的「整体化」得到。

记 $G$ 为 $\mathfrak{g}$ 所对应的伴随群。设我们处于 3.1.2 的情形。于是有 $X$ 上的 $B$-丛 $\mathfrak{F}_{B}$、诱导出的 $G$-丛 $\mathfrak{F}_{G}$，以及 $\mathfrak{g}_{\mathfrak{F}}\otimes\omega_{X}$-torsor $\operatorname{Conn}(\mathfrak{F}_{G})$。设 $D$ 是 $X$ 的有限子概形。记 $\operatorname{Conn}_{D}(\mathfrak{F}_{G})$ 为由 $\operatorname{Conn}(\mathfrak{F}_{G})$ 诱导出的 $\mathfrak{g}_{\mathfrak{F}}\otimes\omega_{X}(D)$-torsor；于是 $\operatorname{Conn}_{D}(\mathfrak{F}_{G})$ 的截面是带 $(\le D)$-奇点的联络。与 3.1.2 一样定义 $c:\operatorname{Conn}_{D}(\mathfrak{F}_{G})\to(\mathfrak{g}/\mathfrak{b})_{\mathfrak{F}}\otimes\omega_{X}(D)$。$X$ 上 $(\le D)$-奇异 $\mathfrak{g}$-oper 的概念如下定义：把定义 3.1.3 中的 $\operatorname{Conn}$ 换成 $\operatorname{Conn}_{D}$，$\omega_{X}$ 换成 $\omega_{X}(D)$。

若 $X$ 完备，则 $X$ 上的 $(\le D)$-奇异 $\mathfrak{g}$-oper 构成概形。与 3.8.9 一样可以证明，从该概形到 $\operatorname{Op}_{\mathfrak{g}}(X\setminus D)$ 的自然态射是闭嵌入，且其像等于 $\operatorname{Op}_{\mathfrak{g},D}(X)$。于是可把 $\operatorname{Op}_{\mathfrak{g},D}(X)$ 看作 $X$ 上 $(\le D)$-奇异 $\mathfrak{g}$-oper 的模概形。

若 $D\subset D'$，则 $\operatorname{Op}_{\mathfrak{g},D}(X)\subset\operatorname{Op}_{\mathfrak{g},D'}(X)$。所以我们应当有自然的办法，从 $(\le D)$-奇异 $\mathfrak{g}$-oper $(\mathfrak{F}_{B},\nabla)$ 构造 $(\le D')$-奇异 $\mathfrak{g}$-oper $(\mathfrak{F}'_{B},\nabla')$。当然 $(\mathfrak{F}'_{B},\nabla')$ 应配有同构 $\alpha:(\mathfrak{F}'_{B},\nabla')|_{X\setminus\Delta}\xrightarrow{\sim}(\mathfrak{F}_{B},\nabla)|_{X\setminus\Delta}$，其中 $\Delta\subset X$ 是把 $D$、$D'$、$\Delta$ 看作有效除子时满足 $D'=D+\Delta$ 的有限子概形。联络 $\nabla'$ 由 $\nabla$ 与 $\alpha$ 重构，而 $(\mathfrak{F}'_{B},\alpha)$ 由下述性质定义（参看 3.8.10）：若 $x\in\Delta$，$f$ 是 $\Delta$ 在 $x$ 处的局部方程，$s$ 是 $\mathfrak{F}_{B}$ 在 $x$ 处的局部截面，则存在 $\mathfrak{F}'_{B}$ 在 $x$ 处的局部截面 $s'$，使得 $\alpha(s')=\lambda(f)s$，其中 $\lambda:\mathbb{G}_{m}\to H$ 是对应于 $\check\rho$ 的态射。
